\section{Conclusion}
\label{sec:conclusion}

This study shows that interoperability of chemical substance data in European regulatory frameworks is fundamentally limited by the lack of a shared and formalised notion of chemical identity. Analysis of 18 regulatory datasets reveals that a substantial fraction of substances cannot be unambiguously linked to a molecular structure, and that commonly used identifiers such as CAS numbers are insufficient for reliable cross-dataset integration.

These limitations directly affect semantic interoperability. Many regulatory definitions lack explicit structure-based identifiers and formal semantics, leaving relationships such as equivalence, inclusion, or overlap implicit. This is particularly evident for substance groups and analytical parameters, which often represent sets or transformations of underlying chemical entities. Without an explicit formalisation of these relationships, interoperability and automated reasoning remain severely constrained.

The problem is amplified by the heterogeneity and scale of regulatory data, where multiple partially overlapping substance lists are defined in different legal and domain-specific contexts. In this setting, both manual curation and data-driven approaches face intrinsic scalability limits, making retrospective harmonisation insufficient.

Our results indicate that interoperability cannot be achieved through post hoc technical solutions alone, including semantic embedding approaches. Instead, it requires changes at the level of data modelling and regulation. The use of structure-based identifiers (e.g. InChIKeys), combined with formal ontologies and linked data principles, is necessary to enable consistent integration when applied at source.

We therefore argue that achieving FAIR and interoperable chemical data requires a shift towards embedding these principles directly into regulatory data standards and legislative processes. Without such changes, fragmentation between regulatory and scientific data infrastructures will persist.

Future work should focus on formal representations for mixtures and substance groups, and on translating these into operational standards for regulatory data exchange.